\section{实现}

\subsection{技术栈}

EdgeSched基于Python 3.10实现。核心框架：FastAPI 0.110+提供WebSocket API，Uvicorn作为ASGI服务器，Pydantic 2.6+用于数据验证。机器学习推理：llama-cpp-python 0.2+用于LLM推理（GGUF格式），PyTorch 2.0+配合Diffusers 0.25+用于图像生成（Stable Diffusion）。异步运行时：asyncio实现基于协程的并发，aiohttp作为异步HTTP客户端。硬件配置：NVIDIA GeForce RTX 5070 Laptop（8GB显存），AMD Ryzen 9 8940HX，32GB DDR5，Windows 11。

\subsection{关键模块实现}

\textbf{动态资源迁移}。算法~\ref{alg:migration}展示了核心迁移逻辑。当图像任务到达时，系统保存LLM状态，释放GPU显存，并异步恢复到CPU。热切换机制确保在GPU恢复期间CPU实例继续服务请求。

\begin{algorithm}[t]
\caption{动态GPU-CPU迁移}
\label{alg:migration}
\begin{algorithmic}[1]
\STATE \textbf{function} \textsc{卸载到CPU}()
\STATE \quad 保存KV缓存状态
\STATE \quad 释放GPU显存
\STATE \quad 异步加载CPU后端
\STATE \textbf{end function}
\STATE
\STATE \textbf{function} \textsc{恢复到GPU}()
\STATE \quad 后台预加载GPU实例
\STATE \quad 等待GPU就绪
\STATE \quad 原子切换：CPU $\rightarrow$ GPU
\STATE \textbf{end function}
\end{algorithmic}
\end{algorithm}

\textbf{任务队列管理}。图像队列使用有界容量防止资源耗尽。队列满时立即拒绝新任务并发送背压信号。

\textbf{状态持久化}。StateBackplane使用mmap实现零拷贝I/O。任务状态通过struct.pack序列化后写入内存映射文件，实现亚毫秒级恢复。

\subsection{部署配置}

\textbf{LLM配置}：
\begin{itemize}
\item 模型：Qwen2.5-7B-Instruct（GGUF格式，Q4\_K\_M量化）
\item 模型权重：约4.7GB
\item KV Cache：约1-2GB（上下文长度4096）
\item 总显存占用：约6-7GB
\item GPU层数：默认-1（自动分配），8GB显存环境下通常加载25-33层
\item 首字超时：30秒
\item 最大生成token数：512
\end{itemize}

\textbf{图像生成配置}：
\begin{itemize}
\item 模型：SDXL Turbo 1.0（fp16精度，约6.7GB）
\item 加载策略：按需加载，完成后卸载以释放显存
\item 最大并发数：1
\item 队列容量：8
\item 总槽位数：9（并发 + 排队）
\item 图像尺寸：512$\times$512
\item 推理步数：20
\end{itemize}
